\begin{figure}[H]
\centering
\footnotesize

\begin{tikzpicture}[
  node distance=8mm and 10mm,
  box/.style={draw, rounded corners, align=left, inner sep=6pt, text width=6.2cm},
  smallbox/.style={draw, rounded corners, align=left, inner sep=6pt, text width=6.2cm},
  arrow/.style={-{Stealth[length=2mm]}, thick},
  execbox/.style={draw, rounded corners, align=left, inner sep=6pt, text width=7cm, fill=blue!10}
]

% --- Entrada
\node[box] (cfg) {
  \textbf{Entrada: configuración declarativa (YAML/JSON)}\\
  \vspace{0.15em}
  \textbf{Objetivo:} \textbf{reproducibilidad} y \textbf{automatización} de escenarios federados.
};

% --- Bloques de secciones
\node[smallbox, below=of cfg] (sections) {
  \textbf{Secciones principales del fichero:}
  \begin{itemize}\setlength\itemsep{0.15em}
    \item \textbf{Dataset}: \textit{data\_dir}, modalidades, sujetos, etiquetado.
    \item \textbf{Split}: \textit{train/val/test}, \textit{seed}, \textit{scaler}, estrategia (global/per-subject).
    \item \textbf{Federation}: modo (manual/auto), nº Fog, asignaciones, porcentajes auto.
  \end{itemize}
};

% --- Validaciones
\node[box, below=of sections] (checks) {
  \textbf{Validaciones de consistencia (pre-ejecución)}\\
  \vspace{0.15em}
  Comprobaciones automáticas de \textbf{coherencia}:
  \begin{itemize}\setlength\itemsep{0.15em}
    \item claves obligatorias y tipos (YAML/JSON),
    \item rangos y sumas de splits,
    \item sujetos disponibles y no solapados,
    \item asignaciones Fog/Edge válidas.
  \end{itemize}
};

% --- Generador
\node[box, below=of checks] (gen) {
  \textbf{Generador de manifiesto}\\
  \vspace{0.15em}
  Script: \texttt{prepare\_swell\_federated.py --config <yaml>}\\
  Procesa la configuración y \textbf{materializa} el experimento.
};

% --- Salida manifiesto
\node[box, below=of gen] (manifest) {
  \textbf{Salida: \texttt{manifest.json} + directorio de particiones}\\
  \vspace{0.15em}
  Contiene:
  \begin{itemize}\setlength\itemsep{0.15em}
    \item configuración completa del experimento,
    \item splits globales (train/val/test),
    \item mapeo \textbf{sujetos $\rightarrow$ Fog},
    \item mapeo \textbf{clientes Edge por Fog},
    \item metadatos (nº sujetos, features, muestras).
  \end{itemize}
};

% --- Lanzador de arquitectura
\node[execbox, below=of manifest] (launcher) {
  \textbf{Lanzador: \texttt{run\_architecture\_from\_config.py}}\\
  \vspace{0.15em}
  Opciones:
  \begin{itemize}\setlength\itemsep{0.15em}
    \item \texttt{--config <yaml>} : carga configuración.
    \item \texttt{--manifest <json>} : aplica manifiesto (opcional).
    \item \texttt{--plan-only} : genera y muestra plan de ejecución.
    \item \texttt{--launch} : lanza procesos federados.
    \item \texttt{--dispatch-config} : envía plan por MQTT.
    \item \texttt{--prepare-splits} : materializa particiones si no hay manifiesto.
  \end{itemize}
};

% --- Plan de ejecución
\node[execbox, below=of launcher] (plan) {
  \textbf{Plan de ejecución generado}\\
  \vspace{0.15em}
  Lista de comandos para:
  \begin{itemize}\setlength\itemsep{0.15em}
    \item Servidor central (Flower).
    \item Nodos Fog (brokers MQTT + agregadores).
    \item Clientes Edge (entrenamiento local).
  \end{itemize}
  Incluye variables de entorno (OTEL, MQTT).
};

% --- Ejecución federada
\node[execbox, below=of plan] (exec) {
  \textbf{Ejecución del ciclo federado}\\
  \vspace{0.15em}
  Proceso iterativo:
  \begin{itemize}\setlength\itemsep{0.15em}
    \item \textbf{Rondas}: servidor $\rightarrow$ fog $\rightarrow$ edge (entrenamiento local) $\rightarrow$ fog (agregación) $\rightarrow$ servidor (FedAvg).
    \item Comunicación vía MQTT.
    \item Monitoreo con OpenTelemetry.
  \end{itemize}
};

% --- Resultados
\node[box, below=of exec] (results) {
  \textbf{Resultados y métricas}\\
  \vspace{0.15em}
  Accuracy, precision, recall, F1 en test set. Logs, trazas OTEL, matrices de confusión.
};

% --- Flechas principales
\draw[arrow] (cfg) -- (sections);
\draw[arrow] (sections) -- (checks);
\draw[arrow] (checks) -- node[right]{\textbf{OK}} (gen);
\draw[arrow] (gen) -- (manifest);
\draw[arrow] (manifest) -- (launcher);
\draw[arrow] (launcher) -- (plan);
\draw[arrow] (plan) -- node[right]{\textbf{si --launch}} (exec);
\draw[arrow] (exec) -- (results);

% --- Rama de error
\node[smallbox, right=of checks, text width=5.8cm] (error) {
  \textbf{Si falla una validación:}\\
  se aborta la generación y se reportan
  \textbf{errores accionables} (clave faltante, split inválido, sujeto no existente, etc.).
};
\draw[arrow] (checks) -- node[above]{\textbf{ERROR}} (error);

% --- Rama plan-only
\node[smallbox, right=of plan, text width=5cm] (planonly) {
  \textbf{Modo --plan-only:}\\
  muestra plan en consola, no ejecuta.
};
\draw[arrow] (plan) -- node[above]{\textbf{si --plan-only}} (planonly);

\end{tikzpicture}

\caption{\textbf{Flujo completo HU-059 + ejecución: configuración $\rightarrow$ validación $\rightarrow$ manifiesto $\rightarrow$ planificación $\rightarrow$ lanzamiento $\rightarrow$ ciclo federado.}
Integra la generación del manifiesto reproducible con el script de lanzamiento \texttt{run\_architecture\_from\_config.py}, culminando en la ejecución del aprendizaje federado Edge--Fog--Cloud.}
\label{fig:complete_flow_hu059_execution}
\end{figure}